This chapter describes the two datasets used to train and evaluate
the methods of this thesis. Each dataset has a \emph{symbolic}
variant (per-cell or per-step one-hot encoding) and a
\emph{visual} variant (per-cell or per-step MNIST image). All data
construction is reproducible from the \texttt{mnlearn.data}
library.

\section{Overview}
\label{sec:datasets-overview}

The two datasets play complementary roles. \emph{Sudoku} is the
main task: a $9 \times 9$ constraint-satisfaction problem with a
cyclic graph encoding the row, column and box \emph{all-different}
constraints (Section~\ref{sec:markov-networks}); MAP inference is
NP-hard on this graph, so it is the harder of the two tasks.
\emph{HMC} (Hidden Markov Chain) is a synthetic sequence-labeling
task on a linear chain; MAP inference is exact in
$O(T K^{2})$ time via the Viterbi algorithm
(Section~\ref{subsec:viterbi}). HMC therefore serves as a
\emph{tractable testbed} for the LP-M3N implementation, with the
empirical correctness check carried out in
Chapter~\ref{chap:experiments}.

For each dataset, the symbolic variant supplies a per-cell one-hot
encoding of the input digit, and the visual variant replaces each
non-blank digit by a randomly sampled image of that digit drawn
from MNIST~\cite{LeCun1998}. The MNIST images are partitioned into
\emph{disjoint pools per split}: the training and validation sets
draw from disjoint partitions of the MNIST training images, and the
test set draws exclusively from the MNIST test images. This
guarantees that no MNIST image used at evaluation time has been
seen during training or model selection.

A side-by-side comparison of the two datasets is given in
Section~\ref{sec:datasets-comparison}.